%%Template file
%Heide Brandtstaedter, h.brandtstaedter@htw-berlin.de
%13.7.2010, Berlin. Last update June 2022, Delft.
%Be careful:  Some packages may be outdated.
%Überarbeitung April 25 HB

\usepackage{amsmath, amssymb, amsfonts}   %LATEX-Mathematik-Formatierungen 


\usepackage{url, hyperref}  %für Hyperlinks, Inernet-Adressen im Literaturverzeichnis zitieren
\hypersetup{ colorlinks=true, linkcolor=blue} 

%Sprachpakete deutsch
\usepackage[T1]{fontenc}	%automatische Silbentrennung  
\usepackage[ngerman, english]{babel}	%Neue Rechtschreibung

%Einbindung von Graphiken
\usepackage{graphicx}								
\usepackage{psfrag}					%Text in Graphiken einbinden
\usepackage{pstricks}
\usepackage{pictexwd}
%\usepackage{epic}
\usepackage{rotating}

\usepackage{verbatim}       %Code TExt Umgebung

%?
\usepackage{tikz}
\usetikzlibrary{matrix, fit}
\usetikzlibrary{backgrounds}

\usepackage{empheq}       % Hervorhebung von Gleichungen, Teil vom Bündel »mathtools«

%für Zitierstyles
%\usepackage[authoryear]{natbib}

%von Heike
%if you want to put several images into one figure:
\usepackage{subfigure}
\newcommand{\goodgap}{%for subfigure package
\hspace{\subfigtopskip}%
\hspace{\subfigbottomskip}}

% Formatierung fuer HTW Thesis
%\documentclass[11pt,a4paper,oneside]{scrartcl}
\usepackage[left=2.5cm,right=2.5cm,top=2.0cm,bottom=2cm]{geometry}%Einstellungen der Seitenränder		
\usepackage{setspace}				%Zeilenabstände


%PGF Plots!
\usepackage{pgfplots}
\pgfplotsset{compat=1.14}
\usepgfplotslibrary{fillbetween,groupplots,units,colorbrewer,patchplots}


%THESIS ERGAENZUNG (OL)
%Paket für erweiterete Tabellenumgebung TabularX
\usepackage{tabularx}													
\usepackage{multirow}%verbundene Zellen und Spalten
\usepackage{booktabs}%verbesserte Optik von Tabellen (midrule, toprule, etc)

% Keine Indents (Einzüge) global:
%\setlength{\parindent}{0pt}

% Länge des automatischen Einzugs bei neuen Absätzen 
%\parindent 0cm																							
% Zeilenabstand
%\onehalfspacing		% 1,5-fachen Zeilenabstand
\addtolength{\footskip}{-6mm}

%Kopf- und Fusszeilen
\usepackage[headsepline]{scrlayer-scrpage}